\begin{abstract}
   Makalah ini membahas peran Matriks Gram dan Convolutional Neural Network (CNN) dalam algoritma \textit{Neural Style Transfer} (NST) untuk sintesis citra artistik. Citra digital terlebih dahulu direpresentasikan sebagai tensor berdimensi tinggi sehingga setiap gambar dapat dipandang sebagai elemen ruang vektor, kemudian dipetakan oleh CNN (VGG-19) menjadi \textit{feature map} pada berbagai kedalaman layer yang membedakan informasi konten dan gaya. Konten direpresentasikan oleh matriks fitur $F_c^l$ dan $F_g^l$ pada layer tertentu, sedangkan gaya dimodelkan melalui Matriks Gram $G_s^l$ dan $G_g^l$ yang menangkap korelasi antar kanal fitur dan bersifat simetris serta \textit{positive semi-definite}.
\end{abstract}

\begin{IEEEkeywords}
    neural style transfer, citra, gaya, matriks gram, aljabar linier, CNN VGG-19
\end{IEEEkeywords}
